\documentclass[10pt]{article}

\usepackage[utf8]{inputenc}

\usepackage[T1]{fontenc}

\usepackage{amsmath}

\usepackage{amsfonts}

\usepackage{amssymb}

\usepackage[version=4]{mhchem}

\usepackage{stmaryrd}

\usepackage{bbold}



\begin{document}

построены методом комплексного ростка для уравнений Шредингера (см. [1]), Клейна-Гордона, псевдодифференциального уравнения типа Шредингера для релятивистской частицы (см. [2]), уравнения Дирака с аномальным взаимодействием Паули (см. [3]) (см. Дирака уравнение; асимптотические решения). В случае уравнения Шредингера со стационарным скалярным потенциалом решения подобного типа впервые построены в [4]. Для произвольной квантовой системы, имеющей классич. аналог, К. т.-К. с. реализуют предельный переход при $h \rightarrow 0$ из квантовой теории в классическую (в духе подхода Эренфеста в нерелятивистской квантовой механике (см. [5])): средние квантово-механические для полного набора наблюдаемых, рассчитанные по К. т.-к.с., в пределе при $h \rightarrow 0$ являются решениями классич. уравнений движения.



Лит.: [1] Багров В. Г., Белов В. В., Тернов И. М., «Теоретич. и математич. физика», 1982, т. 50, № 3, с. $390-96$; [2] Б а гро в В.Г., Белов В. В., «Известия ВУЗов. Физика», 1982, т. 25 , № 4, с. 48-50; [3] Бе л о в В. В., Масло в В. П., “Доклады АН CCCP», 1989, т. 305, № 3, с. 574-80; [4] Баби ч В. М., Дан ило в Ю. П., «Записки научных семинаров ЛОМИ АН СССР», 1969, т. 15, в. 2, c. 47-65; [5] Ehrenfest P., «Z. Phys.», 1927, Bd 45,

S. $455-57$. Белoв.



КВАЗИКЛАССИЧЕСКИЙ ОПЕРАТОР ПлОтности см. P-представление.



КВАЗИКЛАССИЧЕСКОЕ ПРИБЛИХЕНИЕ В И З Л Ч е н и и - приближенные по $h \rightarrow 0$ методы расчета в представлении Фарри характеристик спонтанного излучения электромагнитных волн заряженными частицами, движущимися в произвольном внешнем поле. Эти методы основаны на учете двух типов квантовых поправок к классич. излучению, связанных с квантованием самого движения частицы в этом поле и с эффектом отдачи частицы при излучении (см. [1]). Соответствующие характерные параметры разложения в типичной ситуации (для магнитного поля) $h \omega_{0} / E$ и $h \omega / E$, где $\omega_{0}-$ частота вращения, $\boldsymbol{E}-$ энергия электрона, $\omega-$ частота излученного фотона. В задаче об излучении частицы в области высоких энергий $\gamma=\frac{E}{m c^{2}} \rightarrow \infty$, когда оба типа поправок могут быть разделены, причем квантовые поправки первого типа оказываются малыми по сравнению с эффектом отдачи (см. [2]); основным является операторный квазиклассический с релятивистской точностью метод (см. [3]). В этом методе не требуется знания явного вида решений волновых уравнений, и расчет характеристик процесса излучения проводится в пренебрежении коммутаторами (гейзенберговских) операторов координат и импульсов электрона (величины $\sim h \omega_{0} / E$ ) (подробнее см. [4]). Возможна модификация метода, учитывающая усреднение по начальным состояниям частицы. (О случае квантовой бесспиновой частицы в магнитном поле см. [5].)



Общие формулы характеристик процесса излучения в произвольном внешнем поле, не имеющие ограничений снизу на энергию частицы и равномерно пригодные во всей области энергий $\gamma \in[1, \infty)$, дает метод квазиклассич. траекторно-когерентных состояний (см. [6]). Он основан на использовании полного ортонормированного набора динамич. квазиклассич. траекторно-когерентных состояний оператора Дирака с аномальным взаимодействием Паули (см. Квазиклассические траекторно-когерентные состояния).



В рамках этого метода спектрально-угловое распределение вероятностей и излученной энергии с точностью до первой квантовой поправки включительно представляется в виде конкретного функционала от классич. траектории частицы; под классич. траекторией понимается не только решение уравнения Лоренца, но и решение классич. релятивистского уравнения движения спина типа Барамана-Мишеля-Телегди уравнения. При этом учитываются квантовые поправки, связанные как с отдачей электрона, так и с «квантовым характером траектории». В предельном случае ультрарелятивистских частиц $\gamma \rightarrow \infty$ результаты совпадают (в разложении по $h \rightarrow 0$ ) с формулами операторного квазиклассич. метода. Для нерелятивистских частиц метод дает сравнительно простые общие выражения для характеристик спонтанного излучения заряда с учетом его спиновых свойств. Напр., для вероятности переходов $w^{\uparrow \downarrow}$ с переворотом спина



\[

\begin{aligned}

& w^{\uparrow \downarrow}=\frac{e^{2} h}{6 \pi m^{2} c^{5}} \int_{0}^{\infty} \omega^{3}\left|\vec{S}\left(\omega, \xi,-\xi^{\prime}\right)\right|^{2} d \omega, \\

& \vec{S}\left(\omega, \xi, \xi^{\prime}\right)=\int_{-\infty}^{\infty} \vec{S}\left(t, \xi, \xi^{\prime}\right) \exp (i \omega t) d t,

\end{aligned}

\]



где $\vec{S}\left(t, \xi, \xi^{\prime}\right)-$ решение классич. спиновых уравнений, $\xi, \xi^{\prime}= \pm 1$. ных монокристаллах, Новосиб., 1989; [2] S c h w in g e r J., «Proc. Nat. Acad. Sci. USA», 1954, v. 40, p. 132; [3] Бай ер B. Н., Катко в В. М., “Журнал экспериментальной и теоретич. физики», 1968, т. 55, № 4, с. 1542-54; [4] Байер В. Н., Катков В. М., Фади н В.С., Излучение релятивистских электронов, М., 1973; [5] Ахиезер А. Н., Ласкин Н.В., Шульга Н.Ф., «ДокладыАН СССР», 1988, т. 303, № 1, с. 78-81; [6] Багр ов В. Г., Белов В. В., Ма с л о в В. П., там же, 1989, т. 308, № 1, с. 88-91. КВАЗИКОНФОРМНАЯ КРИВАЯ - образ окрУжности при квазиконформном автоморфизме расширенной комплексной плоскости. Жорданова кривая $\gamma \subset \overline{\mathbb{C}}$ является К. К. в том и только в том случае, если существует константа $c>0$ такая, что для любой пары различных точек $z_{1}, z_{2}(\neq \infty)$ на $\gamma$ выполняется неравенство



\[

\min _{j=1,2} \operatorname{diam}\left(\gamma_{j}\right) \leqslant c\left|z_{1}-z_{2}\right|,

\]



где $\gamma_{j}$ - компоненты $\gamma \backslash\left\{z_{1}, z_{2}\right\}$ (геометрич. характеристика К. к., см. [1]). В последнее время этот класс кривых нашел применения в различных областях анализа. Многочисленные приложения понятия К. К. привели к другим эквивалентным определениям и характеристич. свойствам (cM. [2]).



Лит.: [1] Ал ь ф о р с Л., Лекции по квазиконформным отображениям, пер. с англ., M., 1969; [2] Ge hring F., Characteristic properties $\begin{array}{ll}\text { of quasidisks, Montreal, 1982. } & \text { B. В. Горяйнов, Ю. Е. Хохлов. }\end{array}$ КВАЗИКОНФОРМНОЕ ОТОБРАХЕНИЕ - отображение евклидова пространства в себя, при котором бесконечно малые шары переходят в невырождающиеся эллипсоиды. Точнее, сохраняющее ориентацию отображение $f$ области $G$ в область $G^{\prime}\left(G, G^{\prime} \subset \mathbb{R}^{n}\right)$ кваз и кон фо о м но, если конечна величина коэффициента искажения в точке $a$ :



где



\[

Q(f)=\sup _{x \in G} Q(f, x)

\]



\[

Q(f, a)=\lim _{r \rightarrow 0} \frac{\sup _{|x-a|<r}|f(x)-f(a)|}{\inf _{|x-a|<r}|f(x)-f(a)|}

\]



(вследствие чего К. о. называют также о т о б р а же н и я м и с ограниченным искажением). Если $f$ дифференцируемо в точке $a \in G$, то линейное отображение $f^{\prime}(a)$ преобразует шар в эллипсоид, отношение наибольшей полуоси к-рого к наименьшей равно $Q(f, a)<\infty$. Конформное отображение характеризуется условием $Q(f) \equiv 1$, так что К. о. является естественным его обобщением.



Лит.: Математическая энциклопедия, т. 2, М., 1979, кол. 810-13. КВАЗИЛИНЕЙНОЕ ДИФФЕРЕНЦИАЛЬНОЕ ВЫРАXЕНИЕ - сМ. Дифференциальный оператор.



КВАЗИЛОКАЛЬНАЯ НАБЛЮДАЕМАЯ - СМ. Наблюдаемых алгебра.



КВАЗИЛОКАЛЬНЫЙ ОПЕРАТОР - сМ. Пространственного ослабления корреляций принцип.



КВАЗИМАРКОВСКОЕ ПРИБЛИХЕНИЕ В Те ор и И ту р б у ле н т н о с т и - метод конструирования для изотропного спектра энергии $E(k, t)$ из бесконечной цепочки мо-



КВАЗИКЛАССИЧЕСКИЙ 241



16 Математическая физика





\end{document}